\documentclass[11pt,letterpaper]{article}
\usepackage[T1]{fontenc}
\usepackage{lmodern}
\usepackage[margin=1in]{geometry}
\usepackage{amsmath,amssymb,amsthm,mathtools,booktabs,longtable,array}
\usepackage{graphicx,microtype,xurl,enumitem,fancyhdr}
\usepackage[hidelinks,pdftitle={PA-consistency machines: comparison and new reduction experiments},pdfauthor={Research report prepared for Vladimir Reshetnikov}]{hyperref}
\newcommand{\PA}{\mathrm{PA}}
\newcommand{\Con}{\operatorname{Con}}
\newcommand{\Halt}{\operatorname{Halt}}
\newcommand{\file}[1]{\texttt{\nolinkurl{#1}}}
\newtheorem{theorem}{Theorem}[section]
\newtheorem{lemma}[theorem]{Lemma}
\setlength{\emergencystretch}{3em}
\setlist{itemsep=4pt,topsep=6pt}
\pagestyle{fancy}\fancyhf{}
\fancyhead[L]{\small PA-consistency machines}
\fancyhead[R]{\small Comparison and reduction experiments}
\fancyfoot[C]{\thepage}
\setlength{\headheight}{14pt}
\title{PA-consistency machines:\\comparison and new reduction experiments}
\author{Research report prepared for Vladimir Reshetnikov}
\date{November 13, 2026}
\begin{document}
\maketitle
\begin{abstract}
The requested endpoint is a binary, single-tape Turing machine with at
most 195 working states, started on a blank tape, whose nonhalting is
provably equivalent in PA to consistency of ordinary first-order PA.
\textbf{This investigation did not obtain a verified machine meeting
that endpoint.} The smallest established reference table remains the
published 196-state construction. This report and its archive are a
research result, not a claim to have constructed a 195-state machine.

We recheck the supplied 208-state collection, the specified 209-state
predecessor, and PA196 against frozen source identities. The complete
accepted U208 collection has twenty rooted graph-isomorphism classes;
its final deletion is not an additional deletion available in PA196.
New work includes checked combinations of source rewrites, stronger
physical deletion searches, an equivalent split of the first quantifier
head with explicit Hilbert certificates, and two full-schema
countermodels to tempting axiom deletions. We also identify a precise
HA-strength modification of the supplied iPRA core and distinguish the
published one-way internal PA theorem from the requested biconditional.
All computational evidence is classified by its actual verification
scope. No unsuccessful finite search is represented as a lower bound on
all smaller machines.
\end{abstract}
\setcounter{tocdepth}{1}
\tableofcontents
\clearpage

\section{The target and the evidence boundary}

The intended theorem is
\[
 \PA\vdash\bigl(\neg\Halt(M)\leftrightarrow\Con(\PA)\bigr),
 \qquad |Q_M|\le195.
\]
Here $\PA$ has full first-order induction. The machine has alphabet
$\{0,1\}$, a two-way infinite tape initially filled with zeros, one
transition for each working state and read symbol, and one-cell left or
right moves. The distinguished halting control is excluded from $|Q_M|$.
These are the common conventions of the supplied constructions. A larger
alphabet, a nonblank initial program, or a change in how halt is counted
would require a different comparison.

\begin{center}\small
\begin{tabular}{@{}ll@{}}
\toprule
Input & Frozen revision\\
\midrule
ProveIt & \file{dc523e7b06ffb414093c90aa035500ca283e5f75}\\
BusyBeaver, PA196 & \file{407cdc356f6b586635e71ac92be394255f4bfea4}\\
BusyBeaver, requested PA209 & \file{69d98241a721ead20b4f0b17df9062950b40fbfe}\\
\bottomrule
\end{tabular}
\end{center}

The current ProveIt incoming directory includes earlier responses to
this same research problem. They were inspected before the new searches,
and their results are attributed as prior work. In particular, an
alternative 196-state decoder, the shorter-negation experiment,
distributed separation, diagonal vacuity, and the MP-success rewrite
were already present. Replaying them does not make them new discoveries.
The archive preserves the relevant original license and provenance
notices. No source repository was modified or published.

The report uses four distinct kinds of evidence:
\begin{enumerate}
\item Exact finite checks: table sizes, graph isomorphisms, input hashes,
      complete finite startup traces, and finite proof certificates.
\item Symbolic instruction checks: a specified executor checks arbitrary
      register values under its explicit engine contracts. This is
      stronger than testing a few numerical inputs, but the executor and
      its binding still belong to the trusted argument.
\item Mathematical proofs: the split-head equivalence and the
      countermodels below are schema-level arguments. Their accompanying
      finite controls supplement those arguments.
\item Unfinished endpoints: no new end-to-end Lean or Rocq theorem of
      the requested PA biconditional, and no accepted sub-196 table, is
      delivered here.
\end{enumerate}

\input{sections/comparison}

\section{The calculus relevant to the new logical experiments}

For clarity, write $p(a,b)=\binom{a+b+1}{2}+b$. PA196 uses
\[
 I(a,b)=4p(a,b),\quad M(a,b)=4p(a,b)+1,\quad
 U(a,B)=4p(B,a)+2,\quad E(a,b)=4p(a,b)+3.
\]
Code $0$ denotes truth and code $1=M(0,0)$ is the selected contradiction.
The quantifier constructor's argument order is essential. Write
$N(B)=I(B,1)$, $B\wedge C=N(I(B,N(C)))$, and
$\exists x_a B=N(U(a,N(B)))$. An arrow below means $I$.

Apart from the restricted implicational head $\rho_8$, the six supplied
heads are the following families, with parameters made unrestricted by
the command compiler:
\begin{align*}
 & (\exists x_D\,U(D,H))\to H,\\
 & U(D,U(A,H))\to U(A,U(D,H)),\\
 & M(H,D)\to(E(H,A)\to M(A,D)),\\
 & E(H,D)\to(M(A,H)\to M(A,D)),\\
 & \exists x_D\,E(H,D),\\
 & U(D,A)\to(U(D,A\to H)\to U(D,H)).
\end{align*}
The restricted head is
\[
 ((H\to d)\to D)\to((D\to H)\to(A\to H)),\qquad d\le D.
\]
The existing restricted-head derivation supplies classical implicational
reasoning independently of the first quantifier head.

The gate, when $c\ne a,b$, is
\[
 F(a,b,c)=N\bigl(\Delta_c(a,b)\to(V(a,b,c)\to Q(a,b,c))\bigr),
\]
where
\begin{align*}
 \Delta_c(a,b)&=E(a,b)\to U(c,E(a,b)),\\
 V(a,b,c)&=M(a,b)\to U(c,M(a,b)),\\
 P&=c\to M(c,b),\\
 S&=U\bigl(b,(M(a,b)\to P)\wedge(M(c,b)\to c)\bigr),\\
 Q(a,b,c)&=U(c,S\to N(P)).
\end{align*}
When the guard fails it returns truth. Here the natural number $c$
serves both as a formula code and as a variable label. The code-bounding
lemma makes it fresh for the formula it codes. This interface is the
reason a generic slogan about changing variable names is insufficient.
The original construction and its formal proof inventory are given in
\cite{pa196,result,pa209}.

\input{sections/logic}
\input{sections/search}

\section{How to reproduce and interpret the archive}

Start with \file{README.md}. The default command
\begin{quote}\ttfamily
python3 verify.py
\end{quote}
checks the artifact manifest, the complete selected-table comparison,
the new finite Hilbert certificates and countermodel controls, and the
recorded evidence summaries. More expensive source and physical searches
have separate commands, so reproducing the report does not silently
launch a large search. Input snapshots and source identities are retained
alongside the code. The comparison implementation is inherited from the
incoming research and attributed accordingly; the replay binds its
inputs to the current repository snapshot.

The unchanged upstream source/table check was also run successfully:
\file{scripts/check.sh} checked all 392 PA196 rows, the Palomar source
checks, and 41 generated Lean files. Its output is retained in the
archive. This is the default source check, not the optional full Lean
build and axiom audit. The complete upstream Foundation/Mathlib proof
development was not rebuilt for this report.

The literal reference table is included for reproducibility. Every
experimental table is labeled with its actual count and status. A
passing finite startup, a smaller source program, an optimistic grammar
score, or a source-level proof does not by itself certify a smaller
PA-consistency machine. No file in the archive is labeled as an accepted
195-state construction.

\section{Further research questions}

\paragraph{Can the tape marker and source layout be optimized together?}
The return scanner is a concrete place where one state would suffice.
Shortening its zero-run counter changes how temporary path flags are
recognized, not just how clean source paths are recognized. A successful
redesign must constrain intermediate paths or change the flag encoding
and re-establish the engine contracts. The failed direct deletion is
useful because it identifies that interface precisely.

\paragraph{Can an offset simulation enable a deletion?}
Most established folds return to the same head position. A simulation
relation that permits a fixed state-dependent head displacement enlarges
the design space. It must account for writes on both sides, all incoming
paths, startup, and positive progress. The bounded offset screens in
this investigation do not exhaust such relations.

\paragraph{Which logically equivalent axiom presentation minimizes the
compiled decoder?}
The split-head theorem supplies exact finite proof translations, and the
combined diagonal/distributed emitter shortens the macro word. Neither
has yet lowered the validated physical count. The promising objective
is the size of a verified decoder allocation, including available
history tests and marker constraints. Macro-token count alone is an
inadequate objective.

\paragraph{Can the missing internal implication be completed on the
existing exact trace predicate?}
The proof-producing bit-comprehension argument and bounded positive
simulation give a concrete route. Implementing that route once for
PA196 would make the precise calibration reusable for descendants whose
source transformations are formally certified. This would strengthen
the foundation of the optimization work even without a state saving.

\paragraph{Does the native HA-strength core compile competitively?}
The proposed extension of the ConciseIPRA progress rule changes the
induction strength in the required direction. The next useful artifact
would be an explicit low-level interpreter and state budget for that
core, including certificate representation and enumeration. Its compact
abstract proof checker alone cannot answer the state-count question.

\begin{thebibliography}{9}
\bibitem{pa196} James Chang. \emph{Nonhalting and PA-independence for a
196-state Turing machine}. Paper and source at revision
\file{407cdc356f6b586635e71ac92be394255f4bfea4}.
\url{https://github.com/jamesyc/BusyBeaver/tree/407cdc356f6b586635e71ac92be394255f4bfea4}.
\bibitem{result} James Chang. \emph{Formal theorem statements},
\file{docs/result.md}, same PA196 revision.
\url{https://github.com/jamesyc/BusyBeaver/blob/407cdc356f6b586635e71ac92be394255f4bfea4/docs/result.md}.
\bibitem{pa209} James Chang. The requested PA209 snapshot,
\file{69d98241a721ead20b4f0b17df9062950b40fbfe}.
\url{https://github.com/jamesyc/BusyBeaver/tree/69d98241a721ead20b4f0b17df9062950b40fbfe}.
\bibitem{incoming} Vladimir Reshetnikov. ProveIt, incoming research
archives, revision \file{dc523e7b06ffb414093c90aa035500ca283e5f75}.
\url{https://github.com/VladimirReshetnikov/ProveIt/tree/dc523e7b06ffb414093c90aa035500ca283e5f75/docs/incoming}.
\bibitem{comparison} ProveIt. Imported PA196/PA209 comparison and
reduction-search dossier, same revision, including
\file{Computability/BusyBeaver/Research/PA/pa196-vs-u208/COMPARISON.md}.
\url{https://github.com/VladimirReshetnikov/ProveIt/tree/dc523e7b06ffb414093c90aa035500ca283e5f75/Computability/BusyBeaver/Research/PA/pa196-vs-u208}.
\bibitem{ipra} Vladimir Reshetnikov. ProveIt, \emph{ConciseIPRA},
same revision, including the native scalar rules and their interpretation.
\url{https://github.com/VladimirReshetnikov/ProveIt/tree/dc523e7b06ffb414093c90aa035500ca283e5f75/Logic/PrimitiveRecursiveArithmetic/ConciseIPRA}.
\end{thebibliography}
\end{document}
